% Generated draft: requires mathematical and source review.
\documentclass[11pt]{article}
\usepackage[margin=25mm]{geometry}
\usepackage{fontspec}
\setmainfont{Latin Modern Roman}
\usepackage{amsmath,amssymb,mathtools,unicode-math}
\setmathfont{Latin Modern Math}
\usepackage{enumitem,needspace,xurl,hyperref}
\setlength{\parindent}{0pt}
\setlength{\parskip}{5pt}
\setlength{\emergencystretch}{3em}
\newcommand{\N}{\mathbb{N}}
\newcommand{\Z}{\mathbb{Z}}
\newcommand{\Q}{\mathbb{Q}}
\newcommand{\R}{\mathbb{R}}
\newcommand{\C}{\mathbb{C}}
\newcommand{\F}{\mathbb{F}}
\newcommand{\A}{\mathbb{A}}
\newcommand{\PP}{\mathbb{P}}
\newcommand{\E}{\mathbb{E}}
\newcommand{\eps}{\varepsilon}
\newcommand{\dd}{\,\mathrm{d}}
\begin{document}
\section*{UnsolvedMath ID: 30001302}
\textbf{Source problem code:} catalogue entry
\subsection{Definitions and mathematical statement}
For a variety \(X\), the conventional notation \(CH^c(X)\) denotes the group of codimension-\(c\) algebraic cycles modulo rational equivalence. Tensoring with \(\Q\) gives rational coefficients. The supplied record does not define \(r,K,W_r,A,\psi\), the differential forms, or the correspondence action denoted by \(\Pi^*_{\mathrm{dR}}\). Consequently, the following preserves the source conjecture for its intended data, whose hypotheses must be recovered before this becomes a fully specified mathematical statement.

The conjecture asks whether there exist
\[
\Pi^*\in CH^{r+1}(W_r\times A^{r+1})(K)\otimes\Q
\quad\text{and}\quad q\in\Q\setminus\{0\}
\]
such that
\[
\Pi^*_{\mathrm{dR}}\bigl([\omega_A^{r+1}]\bigr)
=q[\omega_\psi^{r+1}]
\]
and
\[
\Pi^*_{\mathrm{dR}}\bigl([\omega_A^j\overline{\omega}_A^{r+1-j}]\bigr)=0
\qquad\text{for every integer }j\text{ with }1\leq j\leq r.
\]
Thus the source's symbol \(\sim\) means equality up to a nonzero rational scalar. The source imposes the vanishing condition only for \(1\leq j\leq r\); no vanishing condition at either endpoint \(j=0\) or \(j=r+1\) is added.

The notation \((K)\) on the Chow group, the allowed values of \(r\), the cohomology groups containing the displayed classes, and the meaning of the powers and conjugation in these expressions are unspecified. These expressions are retained as source notation; no interpretation as particular tensor products, exterior products, or cohomological projections is assumed.
\subsection{Short English statement}
For the geometric and differential-form data intended by the source, the conjecture asks for an algebraic correspondence with rational coefficients that sends the displayed pure differential class to a nonzero rational multiple of the modular differential class and annihilates every displayed mixed class indexed by \(1\leq j\leq r\). The supplied record omits the definitions and hypotheses needed to specify those data and the correspondence action.
\subsection{Sources}
\begin{itemize}
\item Local registry: supplied UnsolvedMath source record, permanent ID \(30001302\), title \emph{Algebraic Cycles Realizing Modular Differential Forms}, source problem code \texttt{catalogue entry}, in the Oberwolfach Reports Open Problems collection. The supplied proposed year is 2009; no individual proposer is specified.
\item Supplied report citation: \emph{Explicit Methods in Number Theory (2010)}, Oberwolfach Reports, pp.\ 1843--1920, DOI \texttt{10.4171/owr/2009/33}. \url{https://doi.org/10.4171/owr/2009/33}. The background describes these reports as documenting mathematical workshops and their open-problem discussions.
\item Supplied literature reference, labelled ``Darmon Oberwolfach report'': \url{https://ems.press/content/serial-article-files/46233?nt=1}.
\item Supplied literature reference, labelled ``Generalized Heegner cycles theorem'': \url{https://doi.org/10.1215/00127094-2142056}.
\item Literature assessment dated exactly \texttt{2026-08-21}: the supplied curation labels the problem partially solved and states that the conjectural algebraic cycle and complex modular parametrisation remain conditional, while Bertolini--Darmon--Prasanna prove a substantial \(p\)-adic Abel--Jacobi and Rankin-series analogue. This is supplied dated curation derived from the OWR material, not an independently verified assessment. The source explicitly requests independent verification of time-sensitive status and cited claims.
\item The registry labels the statement \texttt{corrected\_verified} and describes the algebraic-cycle realization conjecture as complete with normalized cycle, de Rham, and conjugate-form notation. These are supplied registry qualifications, not independent mathematical verification of this draft.
\end{itemize}
\end{document}
